% /chapters/4-group-actions/01-intro.tex

\section{Introduction}\label{sec:intro_group_actions}

This chapter deals with contributions achieved within various areas of graph theory, and Section~\ref{sec:infinitary_amoeba_graphs} develops a descriptive-set-theoretic theory of the \tFer group acting on infinite graphs, completing the thesis-wide thread on Polish group actions.

Section~\ref{sec:border_tracing} reports a peripheral algorithmic contribution; the remaining sections develop the chapter's main theme, the \tFer group on finite and infinite graphs.

The subsequent sections deal with a specific group action defined on finite graphs called the \emph{group of feasible edge replacements} or the \tFer group, for short.
Section~\ref{sec:recursive_constructions} presents the historic motivation for studying this action by defining the \emph{balancing number} of a graph.
This number is not defined for an arbitrary graph, but only for graphs which are \emph{balanceable}, which is a particular extremal or Ramsey-theoretic property.
The same section then presents a family of graphs which admit a rich \tFer group structure, called \emph{global amoeba} graphs.
These graphs are always balanceable and bounding their balancing numbers is an important question in this research field.
More specifically, this section studies an infinite family of finite trees, obtained through recursive constructions, with rich \tFer groups and computes their balancing numbers, thus providing the answer to several open questions.
Section~\ref{sec:amoeba_trees} presents a general study of structural properties of trees which have rich \tFer groups.
While the existing literature only deals with the case when the graph is finite, the remaining sections of the chapter aim to extend the study of the \tFer group to the infinite case, which is the main topic of discussion of Section~\ref{sec:infinitary_amoeba_graphs}.